\begin{landscape}
\begingroup
\fontsize{10}{13}\selectfont
\fancyhead[L]{\small NACHLASS. STÖRUNGEN DER PALLAS}
\renewcommand{\headrulewidth}{0pt}
\par\medskip
\pdfbookmark[2]{Nachlass VI [5]}{pallas-VI-fragment5}
\begin{center}[5.]\end{center}
\par\noindent\hspace*{\parindent}[Die Störungen der halben grossen Axe ergeben sich genau in derselben Weise, wie in der ersten Rechnung. Unser Beispiel stellt sich hier, wie folgt, wenn constans \(=-\tfrac{2}{8}\lambda\dfrac{10000000}{206265}\) ist, also die 7.\ Decimale als Einheit gewählt wird:
\SourceLeaf{556}
\par
\begingroup\small
\noindent\begin{minipage}[c]{.43\textwidth}
\begin{tabular}{@{}lr@{}}
\(b\)\enspace .\,.\,.\,.\,.\,.\,.\,.\,.\,.\,.\,.\,.&\(0{,}95337\)\\
\(b'\)\enspace .\,.\,.\,.\,.\,.\,.\,.\,.\,.\,.\,.\,.&\(0{,}52101\)\\
\multicolumn{2}{r}{\(Q=69^0\,43'\,11''\)}\\
\(\sin Q\)\enspace .\,.\,.\,.\,.\,.\,.\,.\,.&\(9{,}97221\)\\
\(q\)\enspace .\,.\,.\,.\,.\,.\,.\,.\,.\,.\,.\,.\,.&\(0{,}98116\)\\
\(2-5\dfrac{n'}{n}\)\enspace .\,.\,.\,.\,.\,.&\(8{,}74430\)\\
\hline
\(*\)\enspace .\,.\,.\,.\,.\,.\,.\,.\,.\,.\,.\,.\,.&\(2{,}23686\)\\
constans\enspace .\,.\,.\,.\,.\,.\,.&\(1{,}14727_n\)\\
\hline
&\(3{,}38413_n\)
\end{tabular}
\end{minipage}\hfill
\begin{minipage}[c]{.55\textwidth}
also das Störungsglied
\[
-2422\cos(5M'-2M+69^0\,43'\,11'').
\]
\end{minipage}\par
\endgroup
\par\noindent\hspace*{\parindent}Das vollständige Resultat ist zugleich mit den Störungen des ersten Theils der Epoche in der folgenden Notiz gegeben, da die Argumente dort dieselben sind, und zwar ist immer der Sinus des betreffenden Arguments zu nehmen, so dass also z.\ B. das eben berechnete Störungsglied lautet:
\[
+2422\sin(5\text{\mdseries\AstroSymbols ♃}-2\text{\mdseries\AstroSymbols ⚴}+165^0\,35'\,33'').]
\]
\clearpage
\endgroup
\end{landscape}
